\begin{textsample}{NEG Dim 1 – persona\_grok – Score -1.00 – t550\_grok.txt}  \label{ex:f1_neg_028}
Indonesia 's forests are a treasure trove of biodiversity, yet they face immense threats from deforestation and fires that release massive amounts of carbon emissions.
As Annisa Rahmawati, a forests campaigner for Greenpeace Indonesia, points out, this situation mirrors global crises, such as those affecting the Amazon.
Annisa emphasizes the need for action : companies must sever ties with \textbf{suppliers} linked to deforestation, while consumers should push for products free from deforestation in \textbf{commodities} like \textbf{palm} oil, soya, and meat.
She also calls out governments for their role in enabling these issues.
In 2018, Greenpeace 's campaign rallied 1.5 million people against companies like Wilmar and Mondelez.
This pressure resulted in Wilmar adopting No Deforestation, No Peat, No Exploitation ( NDPE ) commitments and issuing a joint statement on \textbf{monitoring}.
Despite these steps, collaborations in 2019 did not succeed in creating a credible \textbf{monitoring} platform, leading Greenpeace to pull out.

% matched lemmas: commodity, monitoring, palm, supplier
\end{textsample}
